
\begin{theorem} \label{thm:onemaxexpo}
	The expected runtime of the $(1+1)$~Opt-IA$_{\textsc{FirstBest}}$ and the $(1+1)$~Opt-IA$_{\textsc{LastBest}}$ using \IPHfcm{}  %with \textcolor{green}{{\expoHD }} 
	with {\linHD } \new{and any $c\leq 1$} for optimising
	\textsc{OneMax} is $\Theta(n\log n)$ with $\tau=\Omega(n^{1+\epsilon})$ for any arbitrary constant $\epsilon>0$.
\end{theorem}
\begin{proof}
	Let $x_t$ be the current solution at the beginning of the iteration $t$. Immediately after the initialisation, the best seen solution is the current individual $x_1$ itself and the mutation potential \new{for both algorithms} is %$M=$ {\expoHD }$=n^{H(x_1,x_1)/n}=n^0=1$. \textcolor{black}{Similarly the linear mutation is also 
		\textcolor{black}{$1$ due \rev{to} $c\cdot H(x_1,x_1)$ being $0$ and stays $1$ since neutral mutations in \textsc{OneMax} require at least two bits to be flipped.} Note that  $x_t \neq x_{t+1}$ if and only if either $f(x_{t+1})\geq f(x_t)$ or $x_{t+1}$ is reinitialised after $x_t$ is removed from the population due to ageing. Thus, \new{for both algorithmic versions}, the mutation operator flips a single bit at every iteration until ageing is triggered for the first time.
		The improvement probability will be at least $ 1/n $ until either $1^n$ or $0^n$ is sampled. 
		Given that the ageing threshold $\tau$  is at least $n^{1+\epsilon}$ for some constant $\epsilon >0 $, 
		the probability that the current solution will not improve $\tau$ times consecutively is at most 
		$(1-\frac{1}{n})^{n^{1+\epsilon}}=e^{-\Omega(n^\epsilon)}$. 
		Hence, the first optimum will be found before ageing is triggered w.o.p. Given that the ageing operator is not triggered, the expected time to find the first optimum can be shown with a fitness level argument. Let $k$ be the number of $0$-bits in $x_t$. The probability that a mutation improves $x_t$ is thus $k/n$ and the expected time for such an improvement is $n/k$. If we sum over all possible $k$, we obtain the expected runtime $\sum\limits_{k=1}^{n}n/k \leq n \ln n + {\color{black}O(n) } $, which is conditional on ageing not being triggered. If ageing triggers, \new{we pessimistically assume that the first instance that ageing is triggered happens at level $n-1$, and we get the maximum mutation potential which} can be at most $n$. Due to FCM the runtime would be at most $O(n^2 \log n)$ following the same line of argument while considering $n$ wasted fitness function evaluations for each mutation. %Since the reinitialised solutions are generated uniformly at random, if .  
		Since ageing does not trigger \rev{w.o.p.}, the expected runtime is still $O(n \log{n})$.
		
		To prove the lower bound we make a case distinction with respect to whether the ageing is triggered before the optimum is found or not. If the ageing mechanism is triggered early, then by our assumption $\tau=\Omega(n^{1+\epsilon})$ the runtime is $\Omega(n \log n)$. For the second case, the operator flips a single bit per mutation. Since the algorithm is initialised with a uniformly random solution, the number of $0$-bits in the initial solution is at least $n/3$ w.o.p. due to Chernoff bound on the binomial distribution. Similar to our argument for the upper bound, the expected time until the number of $0$-bits in the solution decreases from $k$ to $k-1$ is $k/n$ \new{and FCM will stop with probability 1 every time an extra one-bit is gained.}. Therefore, the expected runtime is  at least $(1-2^{-\Omega(n)}) \cdot \sum_{k=1}^{n/3} k/n = \Omega(n \log n)$. 
	\end{proof}
